\section{Further questions}
\subsection{Can the common norm be weakened again}
The new coefficient two is enough because its radius is less than
$3/4$, while the true digit and factorial errors are below $1/32$.
Replacing it by one would lose this direct uniqueness argument:
$1+1/32>1$. Can retained divisibilities or lattice restrictions exclude
the troublesome boundary cases without adding another witness? A useful
answer must cover signed and unit radices before positivity is recovered;
a small numerical search or a proof assuming the desired signs would not
settle that question.

\subsection{Can contraction and alias choices be optimized together}
An internal-edge contraction changes the cross terms and may make a
previously rejected alias identification or square polarization safe.
Conversely, an alias choice can block a useful contraction. Can one certify
a smaller graph by jointly selecting internal tree edges, source-row
orientations and same-logical aliases, while checking every coefficient?
An optimum proved within a finite declared graph family would be a useful
result, but it would not be a global minimum among height-one equations.

\subsection{Can factorial-free arithmetic improve the same frontier}
The incoming factorial-free construction removes factorial authenticity
by separate prime-power and label-product arguments. Its source is not a
coordinate chart of the shared-error source. A meaningful composition
must verify both the reverse CRT construction at all primes and the
source-sign bootstrap, then measure the new literal incidence structure.
Can it lower dimension without losing the sparse or small-height
advantages of the present signed endpoint?

\subsection{What changes can lower the small scalar degrees}
The retained degree-300 theorem applies only to its fixed arithmetic
frame and constant-weight polynomial-center family. A changed budget
can fall outside its hypotheses; a coordinate shear can instead lower
circuit cost while preserving degree. Further progress should distinguish
these mechanisms explicitly. Any proposed carrier or scale change must
retain an unbounded witness proof, especially for sign branches and
factor cancellation. The rejected early-scale guard deletion is not a
usable premise.

\subsection{Transporting the higher degree projections}
The incoming Lucas package reduces its own 248-variable source by acyclic
monic elimination, producing an explicit octic and higher-degree tradeoffs.
Can its valid pivot patterns be transported to the new 237- or 238-variable
sources, after reconciling their changed coordinates and rows? Such a
construction would need new literal restoration certificates, exact degree
and live-variable proofs, and measured coefficient bounds or full expansions.
No such higher-degree endpoint is claimed in this round; a copied pivot
list without those checks would not constitute the required transport.

\subsection{A kernel checked bridge from the literal arrays}
A focused formalization target is the integral restoration of an actual
emitted residual array, followed by the graph-energy zero-set theorem and
the existing arithmetic calibration. A transitive proof-assistant build
binding the exact array to finite syntactic refutation would strengthen
the present ordinary proof and independent finite checks. Included theorem
names, checksums and successful Python execution do not already provide
that stronger artifact.
